\chapter{Raw Systems Programming \& Zero-Copy Memory (\texttt{std:buffer})}

\section{The Physical Reality of Memory}

\begin{kidbox}{The Street with Billions of Tiny Houses}
Imagine you are walking down the longest street in the world.
On both sides of the street, there are billions of tiny houses, side by side.
Every single house has a street number painted on the front door: House \#0, House \#1, House \#2, all the way to House \#16,000,000,000!
Inside every house, there is a tiny shelf that can hold exactly one byte (a number between 0 and 255).
In computers, this street is called \textbf{RAM} (Random Access Memory), and the house numbers are called \textbf{Memory Addresses}!
\end{kidbox}

\section{Endianness: Reading Backwards or Forwards?}

\begin{kidbox}{The Egg-Breaking War}
In the famous story of \textit{Gulliver's Travels}, two kingdoms went to war over whether soft-boiled breakfast eggs should be cracked open from the Big End or the Little End!
Computers had the exact same argument:
\begin{itemize}
    \item \textbf{Big-Endian}: Writing numbers the way humans write them: most significant byte first. The internet network protocol uses Big-Endian!
    \item \textbf{Little-Endian}: Writing numbers least significant byte first (backwards!). Intel and ARM computer chips use Little-Endian because it makes math circuits faster!
\end{itemize}
Shae's \texttt{ByteBuffer} lets you read and write both with complete precision!
\end{kidbox}

\section{Zero-Copy Memory Slicing}

\begin{kidbox}{Looking Through the Magic Window}
Imagine your friend has a giant coloring book with 5,000 pages of drawings.
You want to show your teacher page 42.
Do you spend three hours and fifty dollars making a photocopy of the entire coloring book just to show one page?
No! You just open page 42 and hold it up!
In computer memory, making copies of big chunks of data is what makes computers slow and drains their battery.
Shae's \texttt{ByteBuffer} supports \textbf{Zero-Copy Slicing}: you create a view into the memory without copying a single byte!
\end{kidbox}

\begin{syntaxbox}{Using \texttt{ByteBuffer} in Shae}
\begin{lstlisting}[language=shaelang]
import "std:buffer";

// 1. Allocate a raw 64-byte buffer
let buf = ByteBuffer.allocate(64);

// 2. Write numeric values with explicit endianness
buf.write_u32_le(0, 0x12345678);
buf.write_f64_be(4, 3.1415926535);

// 3. Read back values
let val_u32 = buf.read_u32_le(0);
let val_f64 = buf.read_f64_be(4);

// 4. Zero-copy sub-slice view
let view = buf.slice(0, 16);
println("Slice byte length: " + view.len()); // 16
\end{lstlisting}
\end{syntaxbox}

\section{Binary Struct Packing and Unpacking}

For network packet serialization and hardware protocols, \texttt{std:buffer} includes a high-performance struct packing engine:

\begin{center}\vspace{8pt}
\centering
\begin{tabular}{lll}
\toprule
\textbf{Format Code} & \textbf{C / Native Type} & \textbf{Size in Bytes} \\
\midrule
\texttt{b} / \texttt{B} & Signed / Unsigned 8-bit integer & 1 \\
\texttt{h} / \texttt{H} & Signed / Unsigned 16-bit integer & 2 \\
\texttt{i} / \texttt{I} & Signed / Unsigned 32-bit integer & 4 \\
\texttt{q} / \texttt{Q} & Signed / Unsigned 64-bit integer & 8 \\
\texttt{f} & IEEE-754 32-bit float & 4 \\
\texttt{d} & IEEE-754 64-bit double & 8 \\
\texttt{s} & Fixed-width ASCII / UTF-8 string & Length dependent \\
\bottomrule
\end{tabular}
\par\vspace{4pt}\small\textbf{Table: Shae Binary Format Specifiers}
\end{center}

\begin{lstlisting}[language=shaelang]
// Packing a network packet header: < = Little Endian, I = u32, 2h = two i16s, 4s = 4-byte string
let packet = ByteBuffer.pack("<I2h4s", [1001, -50, 250, "PING"]);
println("Encoded byte size: " + packet.len()); // 12 bytes

// Unpacking back into high-level values
let unpacked = ByteBuffer.unpack("<I2h4s", packet);
println("Unpacked packet: " + unpacked); // [1001, -50, 250, "PING"]
\end{lstlisting}
